\documentclass[12pt]{article}
\usepackage{amsmath, amssymb}
\usepackage{geometry}
\geometry{margin=1in}
\usepackage{hyperref}
\usepackage{booktabs}

\title{Reaction Tube Propulsion System\\
       \large Corrected Mathematical Description (v2)}
\date{}

\begin{document}
\maketitle

\section*{Design Rationale}

The reaction tube is a \textbf{zero-reaction-mass} propulsion concept. All solid
and fluid mass remains permanently aboard the spacecraft; the only substance that
exits the system is radiation (photons). This trades thrust-to-power ratio for
effectively infinite specific impulse, making the concept suited to extremely
long-duration missions where propellant exhaustion would otherwise be the binding
constraint.

The sole source of external thrust is collimated photon emission. The
projectile and fluid subsystems are internal energy-conversion machinery whose
purpose is to excite the fluid medium in a manner that maximises the fraction of
fusion energy converted into directed, collimated radiation.

\section*{1. System Boundary}

\begin{itemize}
  \item \textbf{Mass boundary:} Closed. Projectile mass $m_p$, fluid mass $m_f$, and
        condensed fluid all remain inside the spacecraft at all times.
  \item \textbf{Energy boundary:} Open. Fusion power $P_{\mathrm{fusion}}$ enters;
        directed photon energy $E_{\mathrm{rad}}$ and waste heat exit.
  \item \textbf{Consequence:} Any internal momentum exchange (projectile
        acceleration, fluid collision, recondensation recoil) is an internal force
        and produces zero net thrust by Newton's third law. Only the photon
        momentum flux couples to the spacecraft's external motion.
\end{itemize}

\section*{2. Energy Accounting}

\subsection*{2.1 Energy inputs}

Total power available per second:
\[
  P_{\mathrm{in}} = P_{\mathrm{fusion}} + R_p \tfrac{1}{2} m_p v_p^2
\]
where $R_p$ is the projectile fire rate and $v_p$ is projectile velocity.
Note that the kinetic energy of the projectile loop is itself ultimately sourced
from $P_{\mathrm{fusion}}$; it is listed separately to track the conversion pathway.

\subsection*{2.2 Energy losses}

Two loss channels reduce the energy available for photon emission.

\paragraph{Phase-change absorption.}
The fluid absorbs energy as it transitions from liquid to gas:
\[
  P_{\mathrm{phase}} = \dot{m}_{\mathrm{phase}}\,\Delta H
\]
where $\Delta H$ is the specific enthalpy of vaporisation (J~kg$^{-1}$) and
$\dot{m}_{\mathrm{phase}}$ is the mass flow rate through the phase transition.

\paragraph{Recondensation.}
When gas recondenses to liquid, the latent heat $\Delta H$ is \emph{released},
not consumed. A fraction $\eta_{\mathrm{rec}} \in [0,1]$ of this released energy
is recoverable (returned to $P_{\mathrm{in}}$ via heat exchange or
thermoelectric conversion); the remainder $1 - \eta_{\mathrm{rec}}$ is lost as
waste heat:
\[
  P_{\mathrm{rec,out}} = \eta_{\mathrm{rec}}\,\dot{m}_{\mathrm{phase}}\,\Delta H
  \qquad
  P_{\mathrm{rec,loss}} = (1 - \eta_{\mathrm{rec}})\,\dot{m}_{\mathrm{phase}}\,\Delta H
\]
In steady state, the net phase-change energy cost is therefore:
\[
  P_{\mathrm{phase,net}} = P_{\mathrm{phase}} - P_{\mathrm{rec,out}}
  = (1 - \eta_{\mathrm{rec}})\,\dot{m}_{\mathrm{phase}}\,\Delta H
\]
The original formulation treated condensation as an additional loss $E_{\mathrm{cond}}$;
this is corrected here. The condensation step releases the same latent heat that
vaporisation absorbed; it is a recovery opportunity, not a second loss.

\subsection*{2.3 Energy available for photon thrust}
\[
  \boxed{
  E_{\mathrm{rad}} = C_{\mathrm{rad}}\!\left(
    P_{\mathrm{in}}
    + \eta_{\mathrm{rec}}\,\dot{m}_{\mathrm{phase}}\,\Delta H
    - \dot{m}_{\mathrm{phase}}\,\Delta H
  \right)
  = C_{\mathrm{rad}}\!\left(
    P_{\mathrm{in}}
    - (1-\eta_{\mathrm{rec}})\,\dot{m}_{\mathrm{phase}}\,\Delta H
  \right)
  }
\]
where $C_{\mathrm{rad}} \in [0,1]$ is the photon collimation efficiency — the
fraction of available energy that is emitted as directed, thrust-producing radiation
rather than isotropic waste heat.

\section*{3. Thrust and Dynamics}

\subsection*{3.1 Net thrust}

Since all mass is conserved internally, the \emph{only} thrust term is the photon
momentum flux:
\[
  \boxed{F_{\mathrm{net}} = \frac{E_{\mathrm{rad}}}{c}
  = \frac{C_{\mathrm{rad}}}{c}\!\left(
    P_{\mathrm{in}} - (1-\eta_{\mathrm{rec}})\,\dot{m}_{\mathrm{phase}}\,\Delta H
  \right)}
\]
The mechanical projectile term $F_{\mathrm{proj}} = \eta_{\mathrm{proj}} R_p m_p v_p$
that appeared in the original formulation is \textbf{removed}. It represented an
internal force pair that cancels identically.

\subsection*{3.2 Spacecraft acceleration}

Non-relativistic regime:
\[
  a_{\mathrm{net}} = \frac{F_{\mathrm{net}}}{M}
  \qquad
  v_{\mathrm{ship}}(t) = \int_0^t a_{\mathrm{net}}(t')\,dt'
\]

Relativistic regime (applicable when $v_{\mathrm{ship}}$ is not negligible relative
to $c$):
\[
  F_{\mathrm{net}} = \frac{d}{dt}\!\left(\gamma M v_{\mathrm{ship}}\right),
  \qquad
  \gamma = \frac{1}{\sqrt{1-(v_{\mathrm{ship}}/c)^2}}
\]

\subsection*{3.3 Thrust-to-power ratio}

A photon rocket has a fundamental thrust-to-power ceiling:
\[
  \frac{F}{P} = \frac{1}{c} \approx 3.33 \times 10^{-9}\ \mathrm{N\,W^{-1}}
\]
This is the highest specific impulse achievable but the \emph{lowest} thrust per
watt of any propulsion concept. The reaction tube system achieves a fraction of
this ceiling:
\[
  \frac{F_{\mathrm{net}}}{P_{\mathrm{in}}} =
  \frac{C_{\mathrm{rad}}}{c}
  \left(1 - \frac{(1-\eta_{\mathrm{rec}})\,\dot{m}_{\mathrm{phase}}\,\Delta H}{P_{\mathrm{in}}}\right)
\]
Maximising this ratio requires simultaneously maximising $C_{\mathrm{rad}}$,
maximising $\eta_{\mathrm{rec}}$, and minimising the phase-change losses relative to
$P_{\mathrm{in}}$.

\section*{4. Dimensionless Performance Parameters}

\subsection*{4.1 Energy loss fraction}
\[
  \Phi = \frac{(1-\eta_{\mathrm{rec}})\,\dot{m}_{\mathrm{phase}}\,\Delta H}{P_{\mathrm{in}}}
  \in [0,1]
\]
\begin{itemize}
  \item $\Phi \to 0$: phase-change losses negligible; system approaches ideal
        photon rocket efficiency.
  \item $\Phi \to 1$: nearly all input energy is lost to phase cycling;
        thrust approaches zero.
\end{itemize}

\subsection*{4.2 Thrust regime number}
\[
  \Gamma = \frac{E_{\mathrm{rad}}}{(1-\eta_{\mathrm{rec}})\,\dot{m}_{\mathrm{phase}}\,\Delta H}
  = \frac{C_{\mathrm{rad}}(1-\Phi)}{\Phi}
\]
\begin{itemize}
  \item $\Gamma \gg 1$: photon thrust dominates over losses.
  \item $\Gamma \sim 1$: transition regime; losses and useful output are comparable.
  \item $\Gamma \ll 1$: loss-dominated; system is not viable.
\end{itemize}

\section*{5. Mass and State Evolution}

\subsection*{5.1 Fluid mass}
In steady-state operation (condensation returns all vaporised fluid):
\[
  \frac{dm_f}{dt} = -\dot{m}_{\mathrm{phase}} + \frac{m_{\mathrm{gas}}}{\tau_{\mathrm{cond}}} = 0
  \quad\Longrightarrow\quad
  \dot{m}_{\mathrm{phase}} = \frac{m_{\mathrm{gas}}}{\tau_{\mathrm{cond}}}
\]
where $\tau_{\mathrm{cond}}$ is the mean condensation timescale.

\subsection*{5.2 Projectile dynamics}
Energy transfer to the fluid per collision reduces projectile velocity:
\[
  \frac{dv_p}{dt} = -\frac{\Delta E_{\mathrm{transfer}}}{m_p\,v_p}
\]
The projectile must be re-accelerated by $P_{\mathrm{fusion}}$ to sustain continuous
operation. The re-acceleration power is part of the $R_p \tfrac{1}{2} m_p v_p^2$
term in $P_{\mathrm{in}}$.

\section*{6. Summary of Corrections from v1}

\begin{center}
\begin{tabular}{lll}
\toprule
\textbf{Quantity} & \textbf{v1 (original)} & \textbf{v2 (corrected)} \\
\midrule
Thrust & $F_{\mathrm{proj}} + F_{\mathrm{photon}}$ & $F_{\mathrm{photon}}$ only \\
Condensation & Additional loss $E_{\mathrm{cond}}$ & Partial recovery $\eta_{\mathrm{rec}}\,\Delta H$ \\
Net phase loss & $\dot{m}(\Delta H + E_{\mathrm{cond}})$ & $(1-\eta_{\mathrm{rec}})\,\dot{m}\,\Delta H$ \\
$E_{\mathrm{rad}}$ & $C_{\mathrm{rad}}(P_{\mathrm{in}}-P_{\mathrm{phase}}-P_{\mathrm{cond}})$ & $C_{\mathrm{rad}}(P_{\mathrm{in}}-(1-\eta_{\mathrm{rec}})\dot{m}\Delta H)$ \\
\bottomrule
\end{tabular}
\end{center}

\section*{7. Open Questions}

\begin{enumerate}
  \item \textbf{Recondensation recovery:} What physical mechanism achieves
        $\eta_{\mathrm{rec}} > 0$? Candidates include heat exchangers, thermoelectric
        converters, and thermoacoustic generators. The upper bound on
        $\eta_{\mathrm{rec}}$ is set by the Carnot efficiency of the hot-cold
        temperature differential across the condenser.
  \item \textbf{Collimation efficiency:} What value of $C_{\mathrm{rad}}$ is
        physically achievable? This depends on the emission spectrum of the excited
        fluid and the geometry of the collimation optics.
  \item \textbf{Projectile loop efficiency:} Is kinetic impact a more effective way
        to excite directed photon emission from the fluid than direct laser or
        microwave coupling? This determines whether the projectile subsystem adds
        net value.
  \item \textbf{Fluid selection:} Which working fluid maximises the ratio
        $E_{\mathrm{rad}} / (P_{\mathrm{phase,net}})$? Candidates should have a
        large photon emission cross-section, a low specific enthalpy of vaporisation,
        and a phase-change temperature compatible with the condenser design.
\end{enumerate}


\section*{8. Carnot-Bounded Recovery Analysis and Fluid Selection}

\subsection*{8.1 Thermodynamic framework}

The condenser operates between a hot reservoir at temperature $T_H$ (the
condensing gas) and a cold reservoir at temperature $T_C$ (the spacecraft
radiator, ultimately rejecting heat to deep space). Any real heat-recovery
device obeys the Carnot bound:
\[
  \eta_{\mathrm{rec}} \leq \eta_{\mathrm{Carnot}} = 1 - \frac{T_C}{T_H}
\]
The condensation temperature $T_H$ is a property of the chosen working fluid
(it equals the boiling point at the operating pressure). $T_C$ is a spacecraft
design parameter; a well-engineered deep-space radiator can reach
$T_C \approx 50$--$200$~K.

For a practical recovery device with second-law efficiency $\epsilon \in [0,1]$
relative to the Carnot limit, the actual recovery fraction is:
\[
  \eta_{\mathrm{rec}} = \epsilon \left(1 - \frac{T_C}{T_H}\right)
\]
Thermoelectric converters achieve $\epsilon \approx 0.10$--$0.20$; thermoacoustic
engines $\epsilon \approx 0.25$--$0.40$; recuperative heat exchangers (which
return energy as heat rather than work) can reach $\epsilon \approx 0.70$--$0.90$
but only reduce the \emph{required} vaporisation power, not directly increase
electrical output.

\subsection*{8.2 Net recoverable power}

The condensation step releases latent heat at rate:
\[
  \dot{Q}_{\mathrm{cond}} = \dot{m}_{\mathrm{phase}}\,\Delta H(T_H)
\]
The power actually recovered and returned to $P_{\mathrm{in}}$ is:
\[
  P_{\mathrm{rec}} = \eta_{\mathrm{rec}}\,\dot{m}_{\mathrm{phase}}\,\Delta H(T_H)
  = \epsilon\!\left(1 - \frac{T_C}{T_H}\right)\dot{m}_{\mathrm{phase}}\,\Delta H(T_H)
\]
The irreducible waste power radiated to space is:
\[
  P_{\mathrm{waste}} = (1 - \eta_{\mathrm{rec}})\,\dot{m}_{\mathrm{phase}}\,\Delta H(T_H)
  + \frac{T_C}{T_H}\,\eta_{\mathrm{rec}}\,\dot{m}_{\mathrm{phase}}\,\Delta H(T_H)
\]
where the second term is the unavoidable heat rejected to the cold reservoir
even in an ideal recovery cycle.

\subsection*{8.3 Updated thrust equation with explicit Carnot terms}

Substituting $\eta_{\mathrm{rec}} = \epsilon(1 - T_C/T_H)$ into the thrust
equation from Section~3.1:
\[
  \boxed{
  F_{\mathrm{net}} = \frac{C_{\mathrm{rad}}}{c}
  \left[
    P_{\mathrm{in}}
    - \left(1 - \epsilon\!\left(1 - \frac{T_C}{T_H}\right)\right)
      \dot{m}_{\mathrm{phase}}\,\Delta H(T_H)
  \right]
  }
\]
This makes the dependence on fluid properties and hardware quality explicit.
Three handles are available to the designer: the collimation efficiency
$C_{\mathrm{rad}}$, the device efficiency $\epsilon$, and the fluid selection
(which sets $T_H$ and $\Delta H$).

\subsection*{8.4 Fluid selection as an optimisation problem}

Define the \textbf{specific thrust figure of merit} for the fluid as the
net radiated energy per unit of phase-change mass flow, normalised by
$P_{\mathrm{in}}$:
\[
  \mathcal{F}(T_H) =
  1 -
  \frac{\left[1 - \epsilon\!\left(1 - T_C/T_H\right)\right]\Delta H(T_H)}
       {P_{\mathrm{in}}/\dot{m}_{\mathrm{phase}}}
\]
Maximising $\mathcal{F}$ requires maximising the effective loss reduction
$\eta_{\mathrm{rec}}\,\Delta H(T_H)$ while minimising $\Delta H(T_H)$ itself.
These two objectives are in tension: high $\Delta H$ means more energy to
recover, but also more energy at risk of being wasted if $\eta_{\mathrm{rec}}$
is limited.

The optimum exists at the fluid where the marginal gain from recovery exactly
offsets the marginal cost of vaporisation. Setting
$\partial \mathcal{F}/\partial T_H = 0$, with $\Delta H$ and $T_H$ linked
through the Clausius--Clapeyron relation
$d(\ln P_{\mathrm{sat}})/d(1/T) = -\Delta H / R$, the optimal boiling
point satisfies:
\[
  \frac{d}{dT_H}\!\left[\eta_{\mathrm{rec}}(T_H)\,\Delta H(T_H)\right]
  = \frac{d}{dT_H}\!\left[\Delta H(T_H)\right]
\]
which reduces to:
\[
  \Delta H(T_H)\,\frac{d\eta_{\mathrm{rec}}}{dT_H} = 0
\]
Since $\Delta H > 0$, the optimum occurs where $d\eta_{\mathrm{rec}}/dT_H = 0$,
i.e.\ where the Carnot efficiency is locally insensitive to $T_H$. Given that
$\eta_{\mathrm{Carnot}} = 1 - T_C/T_H$ is monotonically increasing in $T_H$,
this formally pushes the optimum toward the \emph{highest achievable} $T_H$
--- the fluid that remains condensable but has the highest boiling point
consistent with the collimation optics and structural materials.

\subsection*{8.5 Candidate fluid comparison}

Table~\ref{tab:fluids} lists candidate working fluids ranked by the
key figure of merit $\Delta H / T_H$, which is proportional to the entropy
of vaporisation (Trouton's rule predicts this is roughly constant at
$\approx 88$~J~mol$^{-1}$~K$^{-1}$ for many liquids, but metals and
ionic materials deviate strongly).

\begin{table}[h]
\centering
\caption{Candidate working fluids. $T_H$ = normal boiling point,
$\Delta H$ = specific enthalpy of vaporisation.
$\eta_{\mathrm{Carnot}}$ computed at $T_C = 100$~K.
\label{tab:fluids}}
\begin{tabular}{lrrrr}
\toprule
\textbf{Fluid} & $T_H$ (K) & $\Delta H$ (kJ~kg$^{-1}$) & $\eta_{\mathrm{Carnot}}$ & Notes \\
\midrule
Water          & 373  & 2260 & 0.73 & High $\Delta H$; poor at high $T$ \\
Mercury        & 630  & 295  & 0.84 & Low $\Delta H$; moderately high $T_H$ \\
Lithium        & 1615 & 19\,600 & 0.94 & Very high $\Delta H$ and $T_H$ \\
Sodium         & 1156 & 4370 & 0.91 & Good balance; well studied \\
Cesium         & 944  & 537  & 0.89 & Low $\Delta H$; high $T_H$ \\
Tin            & 2875 & 2490 & 0.97 & Extreme $T_H$; materials challenge \\
\bottomrule
\end{tabular}
\end{table}

\noindent
The analysis favours fluids with \emph{high $T_H$} (maximising $\eta_{\mathrm{Carnot}}$)
and \emph{low $\Delta H$} (minimising gross losses). Cesium and mercury stand out
on both counts among accessible materials. Liquid metals also have the advantage of
high thermal conductivity, which improves heat exchanger effectiveness and raises
the practical $\epsilon$. Lithium has the highest Carnot ceiling but its enormous
$\Delta H$ means any recovery shortfall is severely penalised.

\subsection*{8.6 Coupled optimisation summary}

The full system optimisation is a four-variable problem:
\[
  \max_{C_{\mathrm{rad}},\,\epsilon,\,T_H,\,\dot{m}_{\mathrm{phase}}}
  \quad
  F_{\mathrm{net}} = \frac{C_{\mathrm{rad}}}{c}
  \left[P_{\mathrm{in}} -
  \left(1 - \epsilon\!\left(1-\frac{T_C}{T_H}\right)\right)
  \dot{m}_{\mathrm{phase}}\,\Delta H(T_H)\right]
\]
subject to:
\begin{align}
  \dot{m}_{\mathrm{phase}}\,\Delta H(T_H) &\leq P_{\mathrm{in}}
    & &\text{(energy budget)} \notag \\
  \epsilon &\leq \eta_{\mathrm{Carnot}}(T_H,T_C)
    & &\text{(second law)} \notag \\
  T_H &\leq T_{\mathrm{material}}
    & &\text{(structural/optical limits)} \notag \\
  C_{\mathrm{rad}} &\in [0,1]
    & &\text{(collimation bound)} \notag
\end{align}
The binding constraints in practice are likely to be $T_{\mathrm{material}}$
(the temperature rating of the collimation optics and condenser walls) and
$\epsilon$ (the second-law efficiency of the recovery device). These set a
practical ceiling on $\eta_{\mathrm{rec}}$ that is substantially below the
Carnot limit, making fluid selection and $C_{\mathrm{rad}}$ the primary
engineering levers.


\section*{9. Photon Collimation Efficiency: The Key Constraining Factor}

\subsection*{9.1 \'{E}tendue conservation — the fundamental passive limit}

The single most important constraint on $C_{\mathrm{rad}}$ is not an engineering
difficulty but a fundamental theorem of optics, equivalent to Liouville's theorem
in statistical mechanics: \textbf{no passive optical system can increase the
brightness (specific intensity) of a source.}

The \emph{\'{e}tendue} (optical throughput) of any beam is defined as:
\[
  \mathcal{E} = n^2 A\,\Omega
\]
where $n$ is the refractive index, $A$ is the cross-sectional area of the beam,
and $\Omega$ is the solid angle it subtends. For propagation in vacuum ($n=1$),
$\mathcal{E}$ is conserved through any sequence of lenses, mirrors, or apertures.

A hot fluid source of radius $r_s$ emitting thermally radiates into a hemisphere
($\Omega_{\mathrm{source}} = 2\pi$~sr). A parabolic mirror of focal length $f$
and aperture radius $R$ focuses the source and redirects emission into a beam.
By \'{e}tendue conservation, the output beam's divergence half-angle satisfies:
\[
  \theta_{\mathrm{div}} \geq \frac{r_s}{f}
\]
This is not a limit of imperfect optics — it holds for an ideal, aberration-free
parabolic mirror. The source's finite spatial extent is irreducible: each
off-axis point on the source produces a slightly different beam direction, and
these contributions cannot be separated by any passive element.

\subsection*{9.2 Decomposition of $C_{\mathrm{rad}}$}

The overall collimation efficiency decomposes into three independent factors:
\[
  C_{\mathrm{rad}} = \eta_{\mathrm{coll}} \cdot \eta_{\mathrm{col}} \cdot \eta_{\mathrm{spec}}
\]

\paragraph{Geometric collection efficiency $\eta_{\mathrm{coll}}$.}
The fraction of total emitted power intercepted by the collimating optic.
A mirror subtending solid angle $\Omega_m$ at the source captures at most
$\Omega_m / 4\pi$ of the total isotropic emission. A hemispherical mirror
(the physical maximum) gives $\eta_{\mathrm{coll}} \leq 0.5$.

\paragraph{Collimation quality factor $\eta_{\mathrm{col}}$.}
Not all collected photons contribute thrust equally; only the axial component
of their momentum counts. For a beam with half-angle divergence $\theta_{\mathrm{div}}$,
the mean thrust contribution per photon is $\langle \cos\theta \rangle \approx
1 - \theta_{\mathrm{div}}^2/2$ for small angles. From the \'{e}tendue bound:
\[
  \eta_{\mathrm{col}} \approx 1 - \frac{1}{2}\!\left(\frac{r_s}{f}\right)^2
\]
For a compact high-temperature source ($r_s \sim 10$~cm) with a mirror of focal
length $f \sim 1$~m, $\theta_{\mathrm{div}} \approx 0.1$~rad and
$\eta_{\mathrm{col}} \approx 0.995$. For a spatially extended source
($r_s \sim 1$~m, $f \sim 5$~m), $\theta_{\mathrm{div}} \approx 0.2$~rad and
$\eta_{\mathrm{col}} \approx 0.98$. The angular factor is not usually the
dominant loss term; the collection efficiency is.

\paragraph{Spectral transmission $\eta_{\mathrm{spec}}$.}
The fraction of emitted radiation that passes through the collimating optic
without absorption. This depends on the emission spectrum of the fluid and
the reflectivity/transmissivity of the mirror material across that spectrum.
At $T_H = 1000$~K the emission peaks at $\lambda_{\max} \approx 2.9~\mu$m
(mid-infrared), where most metals have high reflectivity ($> 0.95$).
At $T_H = 3000$~K, $\lambda_{\max} \approx 1~\mu$m (near-infrared), where
mirror performance is generally good. Spectral losses are manageable with
the right mirror material; $\eta_{\mathrm{spec}} \approx 0.85$--$0.95$ is
achievable.

\subsection*{9.3 Practical ceiling for the thermal pathway}

Combining the three factors for the thermal emission case:
\[
  C_{\mathrm{rad}}^{\mathrm{thermal}} \leq
  \underbrace{\eta_{\mathrm{coll}}}_{\leq\,0.50}
  \cdot
  \underbrace{\eta_{\mathrm{col}}}_{\approx\,0.98}
  \cdot
  \underbrace{\eta_{\mathrm{spec}}}_{\approx\,0.90}
  \approx 0.44
\]
In practice, a hemisphere-subtending mirror is not physically realisable
alongside a projectile injection system and a fluid recondenser. A realistic
mirror geometry might subtend $\Omega_m / 4\pi \approx 0.3$--$0.4$, giving:
\[
  C_{\mathrm{rad}}^{\mathrm{thermal}} \approx 0.25\text{--}0.38
\]
This is not improvable by better optics. It is the fundamental passive limit
set by the isotropic nature of thermal emission and the geometry of the tube.

\subsection*{9.4 The Wien temperature–wavelength coupling}

The peak emission wavelength follows Wien's displacement law:
\[
  \lambda_{\max} = \frac{b}{T_H}, \qquad b = 2.898 \times 10^{-3}\ \mathrm{m\cdot K}
\]
Higher $T_H$ shifts emission to shorter wavelengths. This matters for collimation
in two ways. First, diffraction — the fundamental limit on beam quality —
scales with $\lambda$: shorter wavelengths allow smaller diffraction-limited
beam divergence for a given mirror aperture. Second, for very high $T_H$
($> 5000$~K), the peak shifts into the visible and UV, where mirror reflectivity
typically degrades and spectral absorption losses increase. The optimal $T_H$
for the thermal pathway therefore has a soft upper bound around $2000$--$4000$~K,
set by mirror material limits and spectral performance.

\subsection*{9.5 The stimulated emission pathway}

Stimulated emission — radiation produced by a population-inverted medium —
is fundamentally different from thermal emission. Photons are emitted coherently,
in phase with the stimulating radiation field, and in the same direction. The
output is a phase-locked beam whose divergence is set by diffraction from the
aperture, not by source extent. The \'{e}tendue limit does not apply because
stimulated emission is not a passive redirection of thermal radiation; it is
an amplification process.

For a gain medium of length $L$ and aperture $D$, the diffraction-limited
divergence is:
\[
  \theta_{\mathrm{diff}} = \frac{\lambda}{D}
\]
For $D = 0.1$~m and $\lambda = 1~\mu$m, $\theta_{\mathrm{diff}} = 10^{-5}$~rad,
giving $\eta_{\mathrm{col}} \approx 1 - 5 \times 10^{-11} \approx 1$. Combined
with $\eta_{\mathrm{coll}} \to 1$ (all amplified photons exit in the beam
direction) and $\eta_{\mathrm{spec}} \to 1$ (no optic to lose energy through):
\[
  C_{\mathrm{rad}}^{\mathrm{stimulated}} \to 1
\]
This is not merely better than the thermal case — it is a qualitatively
different regime. The entire physics of the system changes.

\subsection*{9.6 The critical question: can the projectile create population inversion?}

Whether the reaction tube can access the stimulated emission pathway depends
entirely on whether projectile kinetic impact can produce population inversion
in the working fluid. This is not a question with an obvious answer, and it
represents the most important open physics question in the concept.

Three mechanisms by which kinetic impact might generate population inversion:

\begin{enumerate}
  \item \textbf{Shock excitation.} A supersonic projectile creates a shock wave
        in the fluid. Shock-heated gas can produce non-equilibrium excited
        state populations on timescales shorter than radiative de-excitation.
        This is the operating principle of gasdynamic lasers and pulsed
        shock-tube lasers. The condition is that the shock velocity must be
        sufficient to selectively populate a metastable upper laser level
        before thermalisation distributes energy across all modes.

  \item \textbf{Collisional excitation selectivity.} If the projectile material
        and the fluid are chosen such that collisional energy transfer preferentially
        populates a specific quantum level (as in atomic collision lasers), the
        impact could produce a transient inversion. This requires careful
        matching of projectile material, fluid species, and energy levels —
        essentially the same constraint as designing a chemical laser.

  \item \textbf{Optical pumping via impact luminescence.} The projectile impact
        generates a burst of radiation (impact flash). If that radiation lies
        within the absorption band of a laser-active impurity dissolved in the
        fluid, the flash optically pumps the impurity into an inverted state.
        This decouples the impact physics from the lasing physics, potentially
        relaxing the constraints.
\end{enumerate}

If \emph{none} of these mechanisms operates, the system is confined to the
thermal emission pathway and $C_{\mathrm{rad}} \leq 0.38$. If any of them
works, $C_{\mathrm{rad}}$ could approach unity, increasing thrust by a factor
of 2.5--4$\times$ at the same input power.

\subsection*{9.7 Revised performance bound}

Combining the $C_{\mathrm{rad}}$ analysis with the full thrust equation from
Section~3.1:

\begin{align}
  F_{\mathrm{net}}^{\mathrm{thermal}} &\leq
    \frac{0.38}{c}\!\left[P_{\mathrm{in}}
    - (1-\eta_{\mathrm{rec}})\,\dot{m}_{\mathrm{phase}}\,\Delta H\right] \\[6pt]
  F_{\mathrm{net}}^{\mathrm{stimulated}} &\leq
    \frac{1.00}{c}\!\left[P_{\mathrm{in}}
    - (1-\eta_{\mathrm{rec}})\,\dot{m}_{\mathrm{phase}}\,\Delta H\right]
\end{align}

The ratio $F_{\mathrm{net}}^{\mathrm{stimulated}} / F_{\mathrm{net}}^{\mathrm{thermal}}
\approx 2.6$ represents the maximum gain from solving the population inversion
problem. Whether this gain is achievable via projectile excitation, or whether
it requires a separate optical pump, is the central engineering question for
the reaction tube concept.

\subsection*{9.8 Summary of $C_{\mathrm{rad}}$ bounds}

\begin{center}
\begin{tabular}{lll}
\toprule
\textbf{Regime} & \textbf{$C_{\mathrm{rad}}$ range} & \textbf{Limiting factor} \\
\midrule
Thermal, poor geometry    & $0.05$--$0.15$ & Mirror solid angle \\
Thermal, good geometry    & $0.25$--$0.38$ & \'{E}tendue (fundamental) \\
Stimulated, partial       & $0.5$--$0.8$   & Gain saturation, cavity losses \\
Stimulated, ideal         & $0.9$--$1.0$   & Diffraction only \\
\bottomrule
\end{tabular}
\end{center}


\section*{10. Array Scaling Constraints}

\subsection*{10.1 Motivation}

A single reaction tube operating at a fusion power $P_{\mathrm{fusion}}$ of
practical magnitude (order 1--100\,MW) is physically large. The natural
engineering response is to parallelise: an array of $N$ identical tubes sharing
a common power source and recondensation infrastructure, each contributing
$F_{\mathrm{net}}/N$ of thrust, with the total summing to $F_{\mathrm{net}}$.
This section shows that parallelisation introduces two new dissipation channels
that are not addressed by phononic crystal (PnC) walls, and that dominate the
thermal budget above a critical array size.

\subsection*{10.2 New dissipation channels in an array}

A single tube has one primary loss pathway: radial heat conduction from the
944\,K cesium fluid through the tube wall to the structural cold sink.
PnC engineering suppresses the phononic component of this pathway.

In a parallelised array, two additional channels appear.

\paragraph{Channel 1: Inter-tube radiative cross-coupling.}
Each tube radiates as a greybody at $T_H = 944$\,K. For an isolated tube the
radiation budget is simple: all emitted power either exits as directed thrust
(fraction $C_{\mathrm{rad}}$) or is deposited on the cold radiator. In an
array, each tube also irradiates its neighbours. The power absorbed by tube
$j$ from tube $i$ per unit length is:
\begin{equation}
  \dot{q}_{i \to j} = F_{ij}\,\varepsilon_{\mathrm{outer}}\,\sigma T_H^4 \cdot 2\pi r_o
  \label{eq:cross}
\end{equation}
where $F_{ij}$ is the geometric view factor between tubes $i$ and $j$,
$\varepsilon_{\mathrm{outer}}$ is the emissivity of the outer tube surface,
$r_o$ is the outer tube radius, and $\sigma = 5.67\times10^{-8}$\,W\,m$^{-2}$\,K$^{-4}$.

For two parallel cylinders of radius $r_o$ with centre-to-centre separation
$s$ ($s \geq 2r_o$), the view factor per unit length is:
\begin{equation}
  F_{12} = \frac{1}{\pi}\left[
    \sqrt{\left(\frac{s}{2r_o}\right)^2 - 1}
    + \arcsin\!\left(\frac{2r_o}{s}\right)
    - \frac{s}{2r_o}
  \right]
  \label{eq:viewfactor}
\end{equation}
For a tightly packed array ($s = 2r_o + \delta$ with small gap $\delta$),
$F_{12} \to 0.5$ — each tube intercepts up to half the lateral radiation from
its nearest neighbour. Interior tubes in a large array are surrounded on all
sides; their effective absorbed cross-radiation power per unit length is:
\begin{equation}
  \dot{q}_{\mathrm{cross}} \approx \sum_{j \in \mathrm{nbrs}} F_{ij}\,
  \varepsilon_{\mathrm{outer}}\,\sigma T_H^4 \cdot 2\pi r_o
\end{equation}
For a hexagonally packed array with six nearest neighbours and $F \approx 0.3$,
an interior tube absorbs approximately $0.3 \times 6 = 1.8$ times its own
lateral thermal emission from its neighbours — substantially increasing the
effective thermal load on the fluid. PnC walls are phononic filters and are
optically transparent; they have no effect on this pathway.

\paragraph{Channel 2: Electronic thermal conductivity scaling.}
As established in the PnC analysis, heat transport in liquid cesium is dominated
by the Wiedemann--Franz electronic contribution:
\begin{equation}
  \kappa_e = L_0\,\sigma_e\,T_H
  \label{eq:WF}
\end{equation}
where $L_0 = 2.44\times10^{-8}$\,W\,$\Omega$\,K$^{-2}$ is the Lorenz number
and $\sigma_e$ is the electrical conductivity of liquid cesium
($\sigma_e \approx 1.5\times10^6$\,S\,m$^{-1}$ at 944\,K), giving
$\kappa_e \approx 17$--$19$\,W\,m$^{-1}$\,K$^{-1}$.

The total electronic heat leak from an array of $N$ tubes scales as:
\begin{equation}
  \dot{Q}_e^{\mathrm{array}} = N \cdot \kappa_e \cdot A_{\mathrm{int}} \cdot
  \frac{\Delta T}{w}
  \label{eq:elec_scaling}
\end{equation}
where $A_{\mathrm{int}}$ is the Cs--wall interfacial contact area per tube,
$\Delta T = T_H - T_C = 844$\,K, and $w$ is the wall thickness. The PnC
contribution suppresses the phononic minority ($\kappa_{\mathrm{ph}} \approx
1$--$2$\,W\,m$^{-1}$\,K$^{-1}$) while leaving $\kappa_e$ unaffected. The
ratio of electronic to total conductivity is:
\begin{equation}
  \frac{\kappa_e}{\kappa_e + \kappa_{\mathrm{ph}}} \approx \frac{18}{19} \approx 0.95
\end{equation}
PnC walls therefore address at most 5\% of the total thermal conductance at
the Cs--wall interface. This fraction does not improve with $N$.

\subsection*{10.3 Critical array size}

Define the \textbf{cross-radiation loading factor} $\Lambda$ as the ratio of
power absorbed by an interior tube from its neighbours to its own internally
generated power:
\begin{equation}
  \Lambda = \frac{\dot{q}_{\mathrm{cross}} \cdot l}{P_{\mathrm{tube}}}
  = \frac{n_{\mathrm{nbr}}\,F\,\varepsilon_{\mathrm{outer}}\,\sigma T_H^4
          \cdot 2\pi r_o\,l}{P_{\mathrm{tube}}}
  \label{eq:Lambda}
\end{equation}
where $l$ is the tube length and $n_{\mathrm{nbr}}$ is the number of nearest
neighbours. When $\Lambda \geq 1$, an interior tube absorbs as much energy
from its neighbours as it generates internally — the array enters a
thermally coupled regime in which each tube's operating point depends on all
others. This is the critical array size beyond which per-tube performance
degrades.

For $r_o = 0.05$\,m, $l = 2$\,m, $n_{\mathrm{nbr}} = 6$, $F = 0.3$,
$\varepsilon_{\mathrm{outer}} = 0.8$ (uncoated metal), $T_H = 944$\,K, and
$P_{\mathrm{tube}} = 100$\,kW:
\[
  \dot{q}_{\mathrm{cross}} \approx 6 \times 0.3 \times 0.8 \times
  5.67\!\times\!10^{-8} \times 944^4 \times 2\pi\,(0.05)\,(2)
  \approx 38\,\mathrm{kW}
\]
\[
  \Lambda \approx \frac{38}{100} = 0.38
\]
This is well below unity for 100\,kW tubes — but at $P_{\mathrm{tube}} =
10$\,kW, $\Lambda \approx 3.8$, meaning each interior tube absorbs nearly
four times its own power from neighbours. Smaller tubes in denser arrays
cross the critical threshold rapidly.

\subsection*{10.4 Mitigation strategies in priority order}

Three strategies address the new dissipation channels, ranked by their
leverage.

\paragraph{Priority 1: Selective outer-surface emissivity.}
The cross-radiation channel scales directly with $\varepsilon_{\mathrm{outer}}$.
Replacing $\varepsilon_{\mathrm{outer}} = 0.8$ (oxidised metal) with a
low-emissivity coating on the inter-tube faces reduces $\dot{q}_{\mathrm{cross}}$
proportionally. Polished refractory metals (W, Mo) achieve
$\varepsilon \approx 0.03$--$0.08$ in the near-infrared at 944\,K:
\[
  \dot{q}_{\mathrm{cross}}^{\mathrm{coated}} \approx 0.05 \times
  \dot{q}_{\mathrm{cross}}^{\mathrm{bare}} \quad\Rightarrow\quad
  \Lambda_{\mathrm{coated}} \approx 0.019
\]
This is the highest-leverage single intervention and it has no mass penalty
relative to PnC fabrication.

\paragraph{Priority 2: Geometric shielding.}
Thin radiation shields (multi-layer insulation geometry) inserted in the
inter-tube gaps reduce the effective view factor:
\begin{equation}
  F_{ij}^{\mathrm{shielded}} = \frac{F_{ij}}{1 + (n_s)(1-\varepsilon_s)/\varepsilon_s}
\end{equation}
where $n_s$ is the number of shield layers and $\varepsilon_s$ is their
emissivity. Two layers of polished molybdenum foil ($\varepsilon_s = 0.05$)
reduce the effective view factor by a factor of $\sim 40$.

\paragraph{Priority 3: Electronic conductance engineering at the interface.}
Suppressing $\kappa_e$ at the Cs--wall boundary requires either a wall material
with a large Lorenz-number mismatch relative to cesium, or a thin electrically
resistive interlayer that creates an effective electronic Kapitza resistance.
Candidate interlayer materials are oxides or nitrides with low electronic
conductivity but compatible thermal expansion coefficients with the PnC
substrate. This is less developed than PnC fabrication and is recommended
for later-stage development once the radiative cross-coupling is resolved.

\subsection*{10.5 Revised efficiency multiplier for array systems}

The efficiency multiplier $\chi$ from the PnC paper must be augmented with
an array cross-radiation penalty term:
\begin{equation}
  \chi_{\mathrm{array}} = \chi_{\mathrm{single}} \cdot
  \frac{1}{1 + \Lambda(\varepsilon_{\mathrm{outer}}, r_o, l, n_{\mathrm{nbr}})}
  \label{eq:chi_array}
\end{equation}
where $\chi_{\mathrm{single}}$ is given by Eq.~(3) of the PnC paper and
$\Lambda$ is defined by Eq.~(\ref{eq:Lambda}). For a well-coated array
($\varepsilon_{\mathrm{outer}} = 0.05$) and moderate tube power
($P_{\mathrm{tube}} \geq 100$\,kW):
\[
  \Lambda \approx 0.019 \quad\Rightarrow\quad
  \chi_{\mathrm{array}} \approx 0.98\,\chi_{\mathrm{single}}
\]
The array penalty is negligible when the outer surface emissivity is properly
managed. Without surface treatment ($\varepsilon_{\mathrm{outer}} = 0.8$),
$\chi_{\mathrm{array}}$ degrades to $0.72\,\chi_{\mathrm{single}}$ for the
100\,kW case, and below $0.20\,\chi_{\mathrm{single}}$ for 10\,kW tubes.

\subsection*{10.6 Summary: thermal management priority stack}

\begin{center}
\begin{tabular}{llll}
\toprule
\textbf{Priority} & \textbf{Channel} & \textbf{Intervention} & \textbf{Leverage} \\
\midrule
1 & Cross-radiation     & Low-$\varepsilon$ outer coating  & $20\times$ reduction in $\Lambda$ \\
2 & Cross-radiation     & Inter-tube radiation shields      & $40\times$ reduction in $F_{ij}$ \\
3 & Electronic conduction & Resistive interlayer at Cs--wall & Suppresses $\kappa_e$ \\
4 & Phononic conduction  & PnC wall engineering              & Suppresses $\kappa_{\mathrm{ph}}$ only \\
\bottomrule
\end{tabular}
\end{center}

\noindent
PnC walls remain worthwhile — they address a real loss channel and add
structural function as a high-$Q$ acoustic cavity relevant to the population
inversion question in Section~9. However, in a parallelised array, outer-surface
emissivity management must be treated as the primary thermal design variable.
A PnC wall with a high-emissivity outer surface is substantially less effective
than a plain metal wall with a low-emissivity outer coating.

\end{document}



